\documentclass[tikz, border=5pt]{standalone}
\usepackage{tikz}
\usepackage{helvet}

\input{../tudcolors.tex}
\usetikzlibrary{arrows.meta, decorations.markings}

\definecolor{Mlight}{RGB}{214,230,248}
% \definecolor{Mmid}  {RGB}{114,163,210}
% \definecolor{Mdark}{tudBlue}

\definecolor{Dlight}{RGB}{252,220,195}
% \definecolor{tudOrange}  {RGB}{237,143, 75}
% \definecolor{Ddark} {RGB}{185, 67, 20}

\definecolor{Ulight}{RGB}{204,232,196}
% \definecolor{Umid}  {RGB}{108,181, 90}
% \definecolor{Udark} {RGB}{ 38,130, 55}

\newcommand{\sector}[5]{%
  \fill[#5]
    (#1:#3) arc[start angle=#1, end angle=#2, radius=#3]
    -- (#2:#4) arc[start angle=#2, end angle=#1, radius=#4]
    -- cycle;}

\begin{document}
\begin{tikzpicture}

\def\rA{0}
\def\rB{2.2}   % Theory / Coding
\def\rC{2.7}   % Coding / Application
\def\rD{4.4}   % outer

% Theory (light)
\sector{90}{210}{\rA}{\rB}{tudBlue!40!white}
\sector{210}{330}{\rA}{\rB}{tudOrange!40!white}
\sector{330}{450}{\rA}{\rB}{tudDarkGreen!40!white}
% \sector{0}{90}   {\rA}{\rB}{tudDarkGreen!40!white}

% Coding (mid)
\sector{90}{210}{\rB}{\rC}{tudBlue!80!white}
\sector{210}{330}{\rB}{\rC}{tudOrange!80!white}
\sector{330}{450}{\rB}{\rC}{tudDarkGreen!80!white}
% \sector{0}{90}   {\rB}{\rC}{tudDarkGreen!80!white}

% Application (dark)
\sector{90}{210}{\rC}{\rD}{tudBlue!90!black}
\sector{210}{330}{\rC}{\rD}{tudOrange!90!black}
\sector{330}{450}{\rC}{\rD}{tudDarkGreen!90!black}
% \sector{0}{90}   {\rC}{\rD}{tudDarkGreen!90!black}

\foreach \a in {90, 210, 330}{
  \draw[white, line width=1.8pt] (0,0) -- (\a:\rD);
}
\draw[white, line width=1.0pt] (0,0) circle (\rB);
\draw[white, line width=1.0pt] (0,0) circle (\rC);

\tikzset{outar/.style={draw=gray!55, line width=1.1pt,
  postaction={decorate},
  decoration={markings, mark=at position 1.0 with
    {\arrow[scale=1.2]{Stealth}}}}}

% \foreach \ca in {150, 270, 30}{
  % \draw[outar] (\ca:0.18)         -- (\ca:{\rD-0.18});
  % % \draw[outar] (\ca:{\rB+0.18})   -- (\ca:{\rC-0.18});
  % % \draw[outar] (\ca:{\rC+0.18})   -- (\ca:{\rD-0.18});
% }

\def\rRot{4.9}
\tikzstyle{rotar} = [draw=gray!55, line width=3.2pt,{Stealth}-]

\draw[rotar]  (39:\rRot) arc[start angle=39, end angle=140, radius=\rRot];
\draw[rotar] (159:\rRot) arc[start angle=159, end angle=260, radius=\rRot];
\draw[rotar] (279:\rRot) arc[start angle=279, end angle=380, radius=\rRot];

\def\rLab{5.0}
\node[font=\sffamily\bfseries\LARGE, text=tudBlue!90!black] at (150:\rLab) {M};
\node[font=\sffamily\bfseries\LARGE, text=tudOrange!90!black] at (270:\rLab) {D};
\node[font=\sffamily\bfseries\LARGE, text=tudDarkGreen!80!black] at  (30:\rLab) {U};

\pgfmathsetmacro{\midAB}{(\rB)/2}
\pgfmathsetmacro{\midBC}{(\rB+\rC)/2}
\pgfmathsetmacro{\midCD}{(\rC+\rD)/2}

\node[font=\sffamily\bfseries, text=tudOrange!90!black, rotate=0]
  at (270:\midAB) {Theory};

\node[font=\sffamily\bfseries, text=white, rotate=0]
  at (270:\midBC) {Coding};

\node[font=\sffamily\bfseries, text=white, rotate=0]
  at (270:\midCD) {Engineering};

\end{tikzpicture}
\end{document}
